\section{Decision model and notation}
\subsection{Block graph, periods, and resources}
Let $B$ denote the blocks admitted to a case and let $A\subseteq B\times B$ contain immediate slope-precedence arcs. An arc $(a,b)$ means that block $a$ must be removed no later than block $b$. The relation is directed from prerequisite to dependent block; implementations may store the reverse adjacency for traversal, but the contract and feasibility check must preserve this meaning. Let $T=\{1,\ldots,H\}$ be the planning periods and $R$ the modeled resources. In the current two-resource cases, these resources represent total material movement and processing feed. For each block $b$, $a_{rb}$ is its consumption of resource $r$, and $c_{rt}$ is the period-specific capacity. Units, resource order, and scaling are part of the case instance, rather than assumptions inferred from a chart.

The binary variable $x_{bt}$ is cumulative: it is one when block $b$ has been extracted by the end of period $t$. Its period increment $u_{bt}=x_{bt}-x_{b,t-1}$ identifies extraction during period $t$, with $x_{b0}=0$. A valid schedule has $u_{bt}\in\{0,1\}$, at most one extraction period per block, and nondecreasing cumulative state. Feasibility requires $x_{bt}\leq x_{at}$ for every precedence arc and period, and $\sum_b a_{rb}u_{bt}\leq c_{rt}$ for each resource and period. These checks are performed on the resulting schedule, not inferred from the method label or its construction order. The distinction matters because an ordering can be precedence-safe and still fail a capacity check if period assignment is implemented incorrectly.

The case boundary also records whether a block is admissible to the ultimate pit, the source dimensions, scenario values, sentinel treatment, and any controls. A schedule score is meaningful only when those inputs and the resource units match the score's denominator. Missing values do not become zero-valued blocks by default; invalid sentinels are rejected or handled under the case-specific import rule. The current web workbench applies these constraints to committed research cases. It does not ingest a live mine block model, perform mine-specific reconciliation, or certify that a field plan is operationally executable.

\subsection{Fixed-destination economics and schedule value}
In fixed-destination CPIT, the destination decision is resolved before scheduling. Write $p_b$ for the resulting block value and let $\delta_t$ be the declared discount multiplier for period $t$. The objective is $\max \sum_{b\in B}\sum_{t\in T}\delta_t p_b u_{bt}$. The repository's scenarios record the discount rate and horizon; implementations and reports must use the same period indexing when constructing $\delta_t$. A rate alone is insufficient to reproduce a value if one implementation discounts at the beginning of the period and another at its end. PhaseFlow compares candidates only within the case's declared convention and keeps this convention in the serialized scenario.

The feasible set is defined jointly by precedence, resource capacities, extraction-at-most-once, and the selected block universe. The ultimate-pit closure constrains that universe, but it is not a period schedule. A value-maximizing closure can exceed period capacities, so its objective cannot be presented as an attainable schedule value. Likewise, a sequence of nested shells is an ordering design, not a proof that the periodized sequence is optimal. The pipeline distinguishes schedule constructors, local improvement methods, and upper-bound calculations so their outputs are not conflated.

For a feasible candidate with value $V$ and a valid upper bound $U>0$ for the same formulation, PhaseFlow reports $g=100(U-V)/U$ percent. This is the gap to the available certificate. It does not identify the unknown integer optimum unless the bound is tight. A comparison is therefore paired by formulation, instance, scenario, and objective convention. When those identities differ, the values are not placed in one ranking. The method traces retain the bound route and feasible value so a reader can inspect the denominator instead of relying on a summary label.

\subsection{Destination choice and formulation boundaries}
The destination-choice extension changes the optimization problem because block values, processing loads, and sometimes admissibility depend on the destination assigned to each block. Let $D_b$ denote the destinations allowed for block $b$. A destination-aware method must select or encode a destination and then schedule the resulting extraction decisions while honoring the resource coefficients associated with those choices. Its feasible set is consequently not the fixed-destination CPIT feasible set with a cosmetic change to the displayed grade. PhaseFlow implements separate destination-choice constructors and improvement routines, records their formulation identity, and pairs their candidate value with a destination-choice bound when such a bound is available.

This separation prevents an invalid denominator comparison. A destination-choice plan may make a block profitable by sending it to a different destination, while a CPIT reference fixes the destination before the schedule is formed. Comparing their raw objectives would combine different feasible regions and potentially different resource loads. The product therefore does not score destination-choice outcomes against the fixed-destination CPIT certificate. The same caution applies to published benchmark values: a source objective is usable only after matching the case identity, destination assumptions, discount convention, horizon, and denominator definition.

The current extension does not model a full chain of stockpile inventory, rehandle, quality blending, reclaim decisions, or recourse after uncertain grade is observed. Those mechanisms introduce state and coupling beyond the block-to-destination choice represented here. Their absence is part of the model scope and is stated because users may otherwise infer operational capabilities from the phrase “destination choice.” The browser exposes the declared model and results; it does not make the omitted operational decisions implicitly.

\subsection{Certificates, candidate schedules, and interpretation}
An upper bound and a feasible schedule answer different questions. The schedule is a concrete lower-bound candidate: it proves that at least its objective value is attainable under the encoded constraints. A valid relaxation or bound procedure limits the objective value attainable by any feasible solution in the same model. The difference between them is an optimization gap, not a prediction interval, a measure of geological uncertainty, or a guarantee of realized production. PhaseFlow preserves the method-specific relation between the feasible schedule and the certificate in each trace.

Several bound routes in this study have distinct meanings. Expected extraction times come from a relaxation and support orderings; they are not themselves feasible schedules. Critical-multiplier runs produce resource-specific relaxation information. The implemented two-resource ExTS route takes the minimum of valid single-resource upper bounds and retains a feasible result from the resulting expected-time orderings. A Bienstock--Zuckerberg formulation is an alternative scalable route to the relaxation, not a stronger relaxation than the same LP. These distinctions affect how a result should be read and which value is a legitimate denominator.

Heuristics such as greedy TopoSort, nested-shell sequencing, sliding-window MILP, restricted neighborhood search, and scenario re-planning can improve candidate quality without proving global optimality. Their runtime and solver status are recorded beside objective values. If a local neighborhood limit stops search, the feasible incumbent remains interpretable, but the stop does not imply that the neighborhood is exhausted. If an exact solver reports a node count unavailable under the current trace schema, the count is represented as null, not as a negative sentinel. These rules keep absent evidence distinct from a measured zero and prevent an implementation detail from being mistaken for a mathematical certificate.

\section{Data contracts and provenance}
\subsection{Case index and manifest}
The case-index document names the schema version, the engine version used to produce the inventory, the declared number of cases, and the case records. Its primary role is to make the replay collection enumerable and to let the freshness gate compare the recorded engine version with the product version. A missing index blocks the release checks; a stale engine version is not repaired by changing only the display title. The browser contract test consumes the same index that the Python producer writes, so a case omitted from the index is not silently treated as a complete release artifact.

Each case manifest supplies identity and provenance for one case. Its required groups describe the case identifier and category, real or synthetic status, engine, scenario, instance, artifact location, lane, gate, flags, controls, best result, and scoreboard. This arrangement keeps scenario facts separate from measured outcomes. The manifest also supports an explicit unavailable or unscored state when source data cannot be used. That state carries a reason; an absent licensed input is not converted to a synthetic replacement and is not reported as a zero result. Producer-side schema validation and browser contract tests check that the two consumers agree on the boundary.

\subsection{Schedule trace}
The schedule-trace document is the replay payload. Its required fields identify schema, case, category, title, role, instance, scenario, bound, controls, method rows, and the paired contract. The optional ensemble record carries learned-policy context when the case uses that analysis. Method entries hold the individual candidates and their measured properties, including status, schedule value, runtime, and method-specific gap where available. This supports comparison at the level of the actual case rather than reconstructing results from aggregate prose.

The trace contract is versioned independently from the application. Version 1.3 explicitly allows unavailable solver node counts to be null; negative sentinel values are invalid because they can be mistaken for measurements. Every committed trace is validated against the schema, and TypeScript fixtures check the browser's interpretation of representative contract values. If a trace is missing or malformed, the release check fails before deployment. The browser does not synthesize a missing method result, and a blank plot is not presented as evidence of a zero objective or zero runtime.

\subsection{Learned model and learned-study reports}
The learned-model interchange identifies the prediction target, ordered input features, network layer definition, feature means and scales, and evaluation metrics. For expected-time inference the eleven serialized feature keys are consumed in order and standardized with the stored training values. The architecture includes a base layer description and five separately seeded members; the executable Python and ONNX paths are checked for numerical parity. The inference result estimates a fractional extraction time. It is not a bound, a probability of feasibility, or a calibrated uncertainty interval.

Separate learned-study documents retain training, holdout, guard-validation, parity, failure-mode, and preview-timing evidence. The split is by deposit seed rather than by individual rows, so correlated scenarios from one generated deposit do not appear in both training and evaluation partitions. The report records seeds, sizes, ensemble seeds, epochs, score summaries, and the independent guard results. Those details let a reader distinguish model fitting from selection and final guard evaluation. Missing study files cause the evidence gate to fail; a model weight file by itself does not establish the validity of the reported evaluation.

\subsection{Readiness snapshot and content-addressed lock}
The readiness document records the product version, digest of the contract catalog, and the state of every declared contract. Each contract row distinguishes available data from an unavailable state and retains provenance and a reason. Readiness is recomputed from the catalog and files on disk, then compared with the committed snapshot. It is therefore a freshness statement about the product's declared evidence set, not a general assertion that every scientific claim is correct or that the application is operationally ready for a mine.

The catalog also pins schema versions, owners, example documents, producer and consumer checks, missing-data behavior, model inputs and outputs, and swap-test status. A schema change must be reflected in the catalog and readiness digest. This makes the boundary auditable across the Python bake, stored JSON, and TypeScript replay. A catalog mismatch or stale digest blocks the gate rather than being silently accepted. The snapshot's available state means the listed artifact exists and meets its contract; scientific interpretation still depends on the study design, provenance, and limitations stated in this report.

\section{Software and analysis architecture}
\subsection{Deterministic case pipeline}
The data pipeline turns a declared instance and scenario into a case manifest, a schedule trace, and supporting model or study evidence. The manifest and trace have separate schemas because provenance and result detail evolve for different reasons. The pipeline validates dimensions, units, precedence, sentinel handling, and resource capacities before method results are written. Feasibility checks are applied to constructed schedules, and reported method values are paired with a bound appropriate to their formulation. The committed case index enumerates the traces, while readiness records whether the complete contract set is present and current.

The expected-time experiment has its own training and evaluation artifacts. Feature normalization is learned from training data and serialized with the model. Five fixed initialization seeds define the ensemble members. Training, selection, and independent guard evaluation use disjoint deposit-seed partitions. Model parity is checked between the serialized computation routes, and the reported study tables are derived from committed results rather than typed independently into the manuscript. The artifact builder included with this report verifies the expected trace count, the two 432-row guard sets, and the five-model inference bundle before generating figures and numeric tables. It fails closed when those inputs or expected shapes change.

\subsection{Static replay application and browser calculations}
The browser application is a static consumer of committed case artifacts with selected in-browser calculations for synthetic examples. It is not backed by a live scheduling API: the application repository reserves the API boundary but does not deploy one. Case selection and charting operate on the declared trace and manifest fields. For supported interactive cases, browser-side schedule calculations are checked against Python fixtures so differences in precedence, capacities, or objective values are detectable. The static page can therefore replay the published research bundle without implying that it is connected to a mine's planning system.

The application keeps measured, predicted, and unavailable quantities visually distinct. Method status and result availability determine whether a value can be drawn. Learned expected times are presented as estimates and do not replace exact expected-time calculations or certified upper bounds. Synthetic examples remain labeled as generated; MineLib cases retain their separate source status. The product's data contracts and frontend tests support this distinction, while the user interface remains a view over committed evidence rather than an independent source of scientific results.

\subsection{Reproduction, generated artifacts, and release boundaries}
Reproduction starts from the repository version, declared pipeline settings, case sources, and serialized model artifacts. The synthetic generator is deterministic for a fixed configuration and seed. Rebuilding the case bundle should reproduce its index, manifests, traces, readiness snapshot, and study summaries. External MineLib inputs are not included when redistribution terms do not permit it; users must obtain those sources independently, and cases without usable source data stay unscored. The artifact builder for this report uses pinned plotting dependencies and emits a receipt with input and output hashes and library versions.

The LaTeX documents separate source prose from generated numeric tables and vector figures. Rebuilding figures from the committed traces keeps plotted counts and distributions synchronized with the records. Generated tables preserve the exact reported values while explanatory text supplies the study interpretation and limitations. A stable generation receipt confirms repeatability of those outputs for the current inputs; it does not validate the source measurements, prove external-data access, or make an incomplete manuscript publication-ready. Release checks, independent review, and publication metadata remain separate gates from local compilation and artifact generation.
